% Modernised reading — fonds Alexandre Grothendieck, shelfmark 112, pages 1-4.
%
% This is an INTERPRETATION, not a transcription. Everything the critical
% apparatus carried moves into footnotes. In any dispute of substance the
% transcription governs: whoever cites Grothendieck must cite batch-01.fr.tex.
%
% Pass: Opus 5 (claude-opus-5), 2026-08-16 — first pass, unchecked against the
% pages by a human. The folder was read whole; it holds one batch, pages 1-4.
%
% This is the folder whose transcription carries the heaviest apparatus of the
% eight read in this pass, and that governs what this reading may do. Where the
% page is illegible the reading says so and stops; it does not complete a
% definition from what the surrounding notation makes plausible. Two of the three
% numbered questions are therefore given with a gap in them, which is the honest
% state of the document.
\documentclass[11pt,a4paper]{article}
% Le préambule commun à toutes les transcriptions du fonds Grothendieck.
%
% Il tient en une idée : **l'incertitude se compose, elle ne se résout pas.**
% Une transcription qui lisse un mot douteux a détruit la seule chose pour
% laquelle on la faisait — un lecteur doit pouvoir savoir, à chaque endroit,
% ce qui était sur le feuillet et ce qui a été ajouté. Les six macros qui
% suivent sont donc l'appareil critique, et non de la décoration.
%
% Les mêmes commandes sont reconnues par `scripts/render.mjs`, qui produit la
% vue de lecture du site. Une macro ajoutée ici doit l'être aussi là-bas,
% faute de quoi le rendu échouera bruyamment — ce qui est voulu : un rendu
% silencieusement faux est pire qu'un rendu absent.

\ProvidesPackage{grothendieck}[2026/08/08 Transcriptions du fonds Grothendieck]

\usepackage{amsmath,amssymb,amsthm}
\usepackage{mathrsfs}
\usepackage{stmaryrd}
% Les diagrammes commutatifs sont partout dans ce fonds : c'est la langue
% même de la géométrie algébrique de Grothendieck, et une transcription qui
% les rendrait en prose ne transcrirait rien.
\usepackage{tikz-cd}
% Français par défaut — c'est la langue des feuillets — mais l'anglais est
% chargé aussi : la traduction et le résumé sont des documents anglais, et la
% typographie française (espaces avant les deux-points) y serait une faute.
% Ces documents basculent par \selectlanguage{english} en tête de corps.
\usepackage[english,french]{babel}
\usepackage{fontspec}
\usepackage{geometry}
\geometry{a4paper,margin=25mm,marginparwidth=28mm}
\usepackage{marginnote}
\usepackage{xcolor}
% ulem, not soul. `soul`'s \st and \ul are built on a character-by-character
% scan that XeLaTeX's font machinery breaks: the document fails with an
% undefined control sequence rather than degrading. ulem does the same two jobs
% and is Unicode-engine safe. `normalem` stops it hijacking \emph, which the
% transcriptions use for Grothendieck's own emphasis.
\usepackage[normalem]{ulem}
\usepackage{hyperref}
\hypersetup{colorlinks=true,linkcolor=black,urlcolor=black,pdfborder={0 0 0}}

% --- Métadonnées du lot -----------------------------------------------------
% Reprises telles quelles par le script de rendu, qui les affiche en tête de
% la vue de lecture. Elles fixent ce que le fichier prétend couvrir : sans
% elles, une transcription flotte sans point d'attache dans un fonds de seize
% mille feuillets.

\newcommand{\folder}[1]{\gdef\grothFolder{#1}}
\newcommand{\batch}[1]{\gdef\grothBatch{#1}}
\newcommand{\leaves}[2]{\gdef\grothFirst{#1}\gdef\grothLast{#2}}
\newcommand{\foldertitle}[1]{\gdef\grothFoldertitle{#1}}
\gdef\grothFolder{?}\gdef\grothBatch{?}\gdef\grothFirst{?}\gdef\grothLast{?}\gdef\grothFoldertitle{}

% --- Le résumé --------------------------------------------------------------
%
% Il ouvre la lecture modernisée, sous le titre « Résumé », et il est composé
% en petit. Ce n'est pas un ornement : c'est le seul passage du document qui
% ne soit pas des mathématiques, et un lecteur doit pouvoir le reconnaître
% avant de l'avoir lu — puis le sauter s'il connaît déjà le sujet, ou s'y
% arrêter s'il ne le connaît pas. Le filet qui le suit marque la reprise du
% corps.

\newenvironment{resume}
  {\small\setlength{\parskip}{0.45em}}
  {\par\medskip\hrule\medskip}

% --- Mots-clés ---------------------------------------------------------------
%
% En anglais, en clôture du résumé de la lecture modernisée : le vocabulaire
% moderne sous lequel chercher ce que les feuillets construisent. Le manifeste
% du site (`npm run manifest`) les extrait pour étiqueter la cote entière —
% cette ligne est la source unique des tags, il n'y a pas de fichier à côté.
\newcommand{\keywords}[1]{\par\smallskip\noindent
  {\footnotesize\textsc{Keywords} --- \textit{#1}}\par}

% --- L'appareil critique ----------------------------------------------------

% Le numéro de feuillet, en marge. C'est l'ancre qui permet de régler un
% désaccord sur une formule en regardant *un* feuillet plutôt qu'une chemise
% de six cents. Le site s'en sert aussi pour tourner les pages du fac-similé
% quand on fait défiler la transcription.
\newcommand{\leaf}[1]{%
  \marginnote{\footnotesize\color{gray!70}\textsf{#1}}%
  \label{leaf:#1}\ignorespaces}

% Illisible : on ne devine pas. Un mot inventé qui se lit comme les autres est
% le pire résultat possible.
\newcommand{\ill}{\textcolor{red!60!black}{[\,\ldots\,]}}

% Lecture douteuse : le mot est donné, et signalé comme incertain.
\newcommand{\uncertain}[1]{\uline{#1}}

% Ajout de l'éditeur — une lettre manquante, un mot sous-entendu.
\newcommand{\add}[1]{\textcolor{blue!50!black}{[#1]}}

% Biffé par Grothendieck lui-même. Ce qu'il a rayé fait partie du document :
% une rature dit souvent où la pensée a changé de direction.
\newcommand{\struck}[1]{\sout{#1}}

% Note du transcripteur, sur l'état matériel du feuillet.
\newcommand{\note}[1]{\par\smallskip{\small\color{gray!60!black}\textit{[#1]}}\par\smallskip}

% Note marginale de Grothendieck, distincte du corps du texte.
\newcommand{\marginal}[1]{\par\smallskip\noindent\hspace{1em}%
  {\small\color{gray!60!black}#1}\par\smallskip}

% --- Titre ------------------------------------------------------------------

\AtBeginDocument{%
  \begin{center}
    {\large\bfseries Fonds Alexandre Grothendieck --- cote n\textsuperscript{o}~\grothFolder}\\[2pt]
    {\itshape\grothFoldertitle}\\[4pt]
    {\small feuillets \grothFirst--\grothLast\ (lot \grothBatch)}\\[2pt]
    {\footnotesize\color{gray!60!black}Transcription établie d'après le fac-similé numérisé
     par l'Université de Montpellier.\\
     Les lectures incertaines sont soulignées, les illisibles notées [\,\ldots\,],
     les ajouts entre crochets.}
  \end{center}
  \vspace{1em}}

\endinput


\folder{112}
\pages{1}{4}
\dating{s.d.}
\watermark{Édition de démonstration\\interprétation personnelle de l'œuvre}
\foldertitle{Linéarisation : notes manuscrites (s.d.) — lecture modernisée du dossier entier}

\begin{document}

\section*{Résumé}

\begin{resume}
Quatre feuilles de listing d'ordinateur, écrites en travers du papier, d'une
écriture très petite et très rapide, avec les colonnes imprimées qui
transparaissent sous chaque ligne. Un seul mot en tête, à l'encre :
\emph{Linéarisation}.

Le mot désigne un procédé et non un objet. En voici le principe, dans un cas
familier. Un espace topologique est un objet compliqué et il n'y a aucune raison
qu'on sache le manipuler. Mais on sait lui associer ses groupes d'homologie, qui
sont des objets d'algèbre linéaire --- des groupes abéliens, ou des espaces
vectoriels --- et sur lesquels on sait tout faire. On a linéarisé : troqué un
objet dont on ne sait rien contre un objet dont on sait tout, en acceptant de
perdre au passage une partie de l'information.

Toute la question est alors de savoir ce qu'on a perdu. Et c'est là que ces
pages travaillent. Elles se placent non pas sur les espaces mais sur les
préfaisceaux sur une petite catégorie $A$ --- ce qui est le cadre général où
Grothendieck situe l'homotopie ---, et elles cherchent à isoler les préfaisceaux
pour lesquels la linéarisation ne perd \emph{rien}. Il les appelle les objets
\emph{lisses}. Symétriquement, elles cherchent ceux pour lesquels elle perd
tout : les objets \emph{négligeables}.

Une fois ces deux classes définies, trois questions sont posées, et elles ont
toutes la même forme : est-ce que la restriction aux objets lisses fait de la
linéarisation une équivalence ? Est-ce qu'elle préserve les groupes
d'extensions ? Le programme est celui de la correspondance de Dold--Kan --- qui,
sur les ensembles simpliciaux, identifie exactement les objets simpliciaux
abéliens et les complexes de chaînes --- transporté sur une catégorie $A$
quelconque. La page le dit d'ailleurs sans détour : \emph{OK pour $A =
\Delta$}, et la vraie question est de savoir jusqu'où l'on peut s'éloigner de
$\Delta$.

Ce dossier est le plus difficile à lire du groupe auquel il appartient, et
cette lecture le dit plutôt que de le masquer : plusieurs des définitions qu'il
pose ont, à leur endroit décisif, un mot qui ne se lit pas.

\keywords{linearization, Dold-Kan correspondence, smooth object, derived category of presheaves, chain complex with twisted coefficients}
\end{resume}

\pagerange{1}{2}
\section*{Le cadre et les deux définitions}

Soit $A$ une petite catégorie et $\hat{A}_{\mathrm{ab}}$ la catégorie des
préfaisceaux abéliens sur $A$. On y prend un objet $L$.

Deux constructions sont supposées connues et ne sont pas rappelées sur la page.
La première est l'intégrale $\int_{A} L$, qui est la manière dont ces notes
notent une colimite --- le sens du signe est fixé par le contexte, non par une
définition. La seconde est un produit noté $*$, qui est une convolution : la
page écrit $L^{(a)} * L^{A^{\circ}}$ là où elle vient d'écrire
$\int_{A} L^{(a)}$, ce qui identifie les deux.\footnote{Aucune de ces deux
notations n'est définie dans le dossier. Le rapprochement entre l'intégrale et
la convolution est celui que la page fait elle-même, par une double barre
d'égalité entre les deux lignes ; ce que chacune signifie séparément reste
implicite. On ne les définit pas ici : ce serait leur donner un contenu que le
dossier n'a pas.}

\subsection*{Lisse}

On dit que $L$ est \emph{lisse} si, pour toute flèche $a \to b$ de $A$, la
flèche naturelle
\[
\int_{A_{a}} L_{a} \longrightarrow \int_{A_{b}} L_{b}
\]
est un isomorphisme. Les deux membres sont ensuite réécrits, par une chaîne
d'égalités que la page pose verticalement :
\[
\int_{A_{a}} L_{a} = \int_{A} L^{(a)} = L^{(a)} * L^{A^{\circ}},
\qquad
\int_{A_{b}} L_{b} = \int_{A} L^{(b)} = L^{(b)} * L^{A^{\circ}}.
\]

La condition est donc : la valeur de la convolution ne dépend pas du point où
on la prend.\footnote{Le mot qui qualifie la flèche est écrit très vite et n'est
pas assuré : la transcription le laisse en blanc. Ce pourrait être
« isomorphisme » ou « quasi-isomorphisme », et la différence n'est pas anodine
--- dans un cadre dérivé c'est la seconde qui est la bonne notion. On a retenu
la première parce que c'est celle qui rend la définition la plus faible, donc
la classe des objets lisses la plus large, et parce que rien sur la page
n'impose l'autre ; le choix est le nôtre et il est signalé.} Une lecture
possible, et qui est celle que le mot \emph{lisse} suggère, est que la classe
ainsi définie est celle des préfaisceaux « localement constants » du point de
vue de la linéarisation : ceux dont la linéarisation ne varie pas le long des
flèches de $A$.\footnote{Cette glose est la nôtre. Le mot « lisse » est de lui ;
son rapprochement avec la constance locale ne l'est pas, et le dossier ne le
justifie nulle part.}

\subsection*{Négligeable}

La seconde définition est l'exact complément de la première, et c'est celle qui
souffre le plus de l'état de la page : on dit que $L$ est \emph{négligeable} si
$L^{b}$ vérifie une condition qui n'est pas lisible.\footnote{La ligne est
écrite par-dessus une portion particulièrement chargée du listing imprimé et
deux mots y disparaissent. La ligne suivante --- « c'est-à-dire si » --- reprend
la même chose sous une forme intégrale, mais s'achève elle aussi sur un membre
illisible. La définition est donc, dans cette édition, incomplète des deux
côtés ; on ne la complète pas.}

Ce que la page fournit ensuite est l'objet sur lequel tout repose : on pose
\[
L^{b} = \bigl(a \mapsto L^{(a)} \otimes \ldots \bigr)
\qquad\text{dans } D_{\bullet}(\hat{B}_{\mathrm{ab}}),
\]
et l'on note que c'est un complexe de chaînes à coefficients
tordus.\footnote{Le second facteur du produit tensoriel n'est pas lisible, non
plus que l'adjectif qui qualifie le complexe ; la torsion des coefficients,
elle, est écrite et lisible. C'est le point important : $L^{b}$ n'est pas un
complexe à coefficients constants, et c'est ce qui empêche de traiter la
linéarisation comme une simple homologie.}

On note enfin $\mathrm{Lisse}_{\mathbb{Z}}(A)$ la sous-catégorie pleine de
$\hat{A}_{\mathrm{ab}}$ formée des objets lisses.

\pagerange{2}{2}
\section*{Les trois questions}

\subsection*{(1) Une équivalence ?}

La première question est la principale :
\[
W_{y}^{-1}\,\mathrm{Lisse}_{\mathbb{Z}}(A) \longrightarrow
D^{\bullet}_{lc}(\hat{B}_{\mathrm{ab}})
\]
est-ce une \emph{équivalence de catégories} ?\footnote{Le mot est souligné sur
la page. L'indice $lc$ se lit clairement au feuillet suivant, où il reparaît
deux fois ; il se développe vraisemblablement en « localement constant », ce qui
serait cohérent avec le sens donné plus haut à \emph{lisse}, mais la page ne le
développe jamais.}

Ce que cette flèche fait est ceci : elle prend les objets lisses, inverse une
classe de flèches, et compare le résultat à une catégorie dérivée. C'est
exactement la forme de la correspondance de Dold--Kan, et la page le dit
elle-même en donnant sa propre portée :

\begin{quote}
OK pour $A = \Delta$ ; OK pour $A = \Delta_{/X}$, où $X \in
\mathrm{Ob}\,\hat{\Delta}$ ?
\end{quote}

Le premier cas est affirmé, le second est demandé. Autrement dit : la
correspondance est acquise sur la catégorie des simplexes, et la question est de
savoir si elle survit au passage aux catégories de tranches --- premier pas hors
de $\Delta$, et le plus modeste qu'on puisse faire.\footnote{Une note écrite en
regard, à demi lisible, rappelle que la construction se fait « comme
$M_{\bullet} \mapsto M^{b}$ », c'est-à-dire par le foncteur qui a été défini
plus haut. Un mot de cette note est illisible.}

\subsection*{(2) Les extensions}

Si $L$ et $M$ sont dans $\mathrm{Lisse}_{\mathbb{Z}}(A)$, a-t-on
\[
\mathrm{Ext}^{i}_{\mathbb{Z}_{A}}(L, M) \longrightarrow
\mathrm{Ext}^{i}_{\mathbb{Z}_{B}}(L^{b}, M^{b}) \;?
\]

C'est la question 1 rendue quantitative : une équivalence de catégories induit
des isomorphismes sur les groupes d'extensions, et vérifier ces derniers est le
moyen usuel d'attaquer la première. Deux remarques l'accompagnent entre
crochets. La première demande si cela vaudrait \emph{sans} supposer $L$ et $M$
lisses --- c'est-à-dire si l'hypothèse de lissité est nécessaire ou seulement
commode. La seconde, écrite plus haut, note que
$D_{\bullet}(\mathrm{Ab}) \cong \mathrm{Hot}_{ab}$ lorsque $A$ est
asphérique.\footnote{Cette dernière est le cas dégénéré qui sert de contrôle :
quand $A$ est asphérique, la catégorie dérivée des groupes abéliens est
équivalente à la catégorie homotopique abélienne, et la question 1 y est donc
triviale. Que ce contrôle soit noté juste à côté de la question générale
n'est pas indifférent.}

\subsection*{(3)}

La troisième question commence --- « supposons $L$ lisse » --- et sa suite n'est
pas lisible.\footnote{Elle est écrite en bas de feuillet, sur la portion la plus
sombre du listing. On ne restitue rien.}

\pagerange{3}{3}
\section*{Les trois programmes}

Le troisième feuillet reprend le tout sous une forme plus ordonnée. Il pose
d'abord la correspondance
\[
L_{\bullet} \longmapsto \varphi(L_{\bullet}),
\qquad
H_{i}(L_{\bullet}) \simeq H_{i}\,\varphi(L_{\bullet}),
\]
c'est-à-dire : le foncteur $\varphi$ préserve l'homologie.\footnote{Un
caractère se glisse entre $H_{i}$ et $\varphi$ dans le second membre et n'est
pas identifié ; il pourrait s'agir d'un signe de composition. On le laisse de
côté, ce qui donne l'énoncé le plus simple compatible avec la ligne.} C'est la
condition minimale pour qu'une linéarisation mérite son nom.

Trois programmes sont ensuite listés :

\begin{itemize}
\item[A)] $D^{b}_{lc}(\hat{A}_{k})$, etc.
\item[B)] $D^{-}_{lc}(\hat{A}_{h}) \rightleftarrows D^{-}_{lc}(\hat{B}_{h})$,
avec le mot \emph{équivalence} souligné --- c'est la question 1, reposée avec
des catégories dérivées bornées d'un côté seulement.\footnote{Une note ajoute :
« NB : on a $D^{-}(\hat{A}_{k}) \to D^{-}(\hat{B}_{h})$ », suivie d'un membre
illisible. Les indices $k$ et $h$ alternent d'une ligne à l'autre sans que la
page dise s'ils désignent la même chose ; on les reproduit tels quels plutôt
que de les uniformiser.}
\item[C)] Les objets \emph{lisses} de $\hat{A}_{k}$, dans un cadre annoncé comme
plus général, avec une notation $\hat{A}_{k\text{-lisse}}$ substituée à un
$\mathrm{Lisse}_{k}$ barré, et un $W$ isolé.
\end{itemize}

Ce changement de notation --- un indice à la place d'un nom --- est le seul
signe, dans tout le dossier, d'un travail de mise au propre. Il n'a pas
abouti : le quatrième feuillet reprend les mêmes matières sans les mener plus
loin, et rien n'y forme un énoncé qu'on puisse détacher.\footnote{On ne le
découpe donc pas en propositions qu'il ne porte pas. C'est le seul feuillet
écrit du dossier dont cette lecture ne tire rien.}

\end{document}
